\scanpage{21}
\mspage{\unsure{15}}

me $R_s$ est un tore maximal dans $G_s$, et la formation de $\underline{\operatorname{Norm}}$, $\underline{\operatorname{Centr}}$, $\underline{\operatorname{Norm}}/\underline{\operatorname{Centr}}$ est compatible avec toute extension de la base, en particulier le passage aux fibres, on voit que l'on a
\[
w(s)=w'(s)\quad;
\]
je dis que de plus, pour $t$ voisin de $s$, on aura les inégalités
\[
w(t)\geq w'(t)\quad\text{et}\quad w'(t)\geq w'(s),
\]
ce qui suffira pour établir a). \deleted{La deuxième} Les deux inégalités \deleted{\unsure{résulte}} sont contenues dans les deux lemmes suivants:

\margin{\figure{\begin{tikzpicture}\draw (0,0) .. controls (.25,.65) and (.45,.85) .. (.65,.72);\end{tikzpicture}}}
\uline{Lemme} 2.2. Soient $S$ un préschéma, $W$ un $S$-préschéma en groupes qui est étale, \deleted{\unsure{quasi-fini et}} séparé et quasi-fini sur $S$. Pour tout $s\in S$, soit $f(s)$ la classe de la fibre géométrique de $W$ en $s$, qui est un \deleted{ensemble}\inserted{élément} de l'ensemble ordonné $E$ des classes de groupes finis à isomorphisme près, introduit dans 2.1. Alors la fonction $f:S\longrightarrow E$ est semi-continue inférieurement. Pour qu'elle soit constante au voisinage de $s\in S$, \deleted{(\unsure{ou ce qui revient au même, pour}} il faut et suffit qu'il en soit de même de la fonction $s\rightsquigarrow\operatorname{card}f(s)$, et pour ceci il faut et suffit que $W$ soit fini sur $S$ au dessus d'un voisinage ouvert $U$ de $s$.

Ce résultat est un raffinement, en termes de groupes, de celui invoqué dans la démonstration de \inserted{Exp $\overline{\underline{\mathrm{XI}}}$} 5.10., nous nous bornons à une esquisse de \inserted{la} démonstration (\deleted{de}\inserted{qui est} type \inserted{le} plus standart). On est ramené comme d'habitude au cas $S$ affine noethérien. On voit de suite que la fonction $f$ est constructible (EGA \deleted{\unsure{IV}} $0_{\rm III}$ 9.3.1. et 9.3.2., et sorites de EGA IV 9), et en vertu de EGA $0_{\rm III}$ 9.3.4. on est ramené \inserted{pour la semi-continuité} à prouver que si $t$ est une générisation de $s$, on a $f(t)\geq f(s)$. Cela nous ramène \deleted{au cas} grâce à EGA II 7.1.7. au cas où $S$ est le spectre d'un anneau de valuation discrète, qu'on peut supposer complet à corps résiduel
\mspage{20}
